% Generated draft: requires mathematical and source review.
\documentclass[11pt]{article}
\usepackage[margin=25mm]{geometry}
\usepackage{fontspec}
\setmainfont{Latin Modern Roman}
\usepackage{amsmath,amssymb,mathtools,unicode-math}
\setmathfont{Latin Modern Math}
\usepackage{enumitem,needspace,xurl,hyperref}
\setlength{\parindent}{0pt}
\setlength{\parskip}{5pt}
\setlength{\emergencystretch}{3em}
\newcommand{\N}{\mathbb{N}}
\newcommand{\Z}{\mathbb{Z}}
\newcommand{\Q}{\mathbb{Q}}
\newcommand{\R}{\mathbb{R}}
\newcommand{\C}{\mathbb{C}}
\newcommand{\F}{\mathbb{F}}
\newcommand{\A}{\mathbb{A}}
\newcommand{\PP}{\mathbb{P}}
\newcommand{\E}{\mathbb{E}}
\newcommand{\eps}{\varepsilon}
\newcommand{\dd}{\,\mathrm{d}}
\begin{document}
\section*{UnsolvedMath ID: 30006381}
\textbf{Source problem code:} catalogue entry
\subsection{Definitions and mathematical statement}
Let \(k\) denote the base field and let \(X\) and \(Y\) denote schemes over \(k\), in the setting intended by the supplied report passage. The notation \(\mathrm{DM}(S)\) denotes the stable \(\infty\)-category of étale motives intended there. The supplied excerpt does not specify its coefficients, its precise construction, or the general hypotheses on \(k\) and \(X\). It also does not explicitly define \(\mathrm{DM}(S)^{\prime}\) or the categorical setting for the relative tensor product below. These omissions remain unresolved.

For context, an object \(P\) in a symmetric monoidal category is dualizable if it admits a dual \(P^\vee\), an evaluation \(P^\vee\otimes P\to\mathbf{1}\), and a coevaluation \(\mathbf{1}\to P\otimes P^\vee\) satisfying the two triangle identities. The report calls dualizable objects of \(\mathrm{DM}(S)\) motivic local systems and describes them as central to the conjecture. This description does not by itself specify whether the prime denotes the dualizable subcategory or some completion or enlargement of it.

Let \(p_X:X\times_kY\to X\) and \(p_Y:X\times_kY\to Y\) be the projections. Exterior product is the construction
\[
M\boxtimes N=p_X^*M\otimes p_Y^*N.
\]
\textbf{Conjecture 1, as supplied.} For every pair \(X,Y\) in the report's intended class of \(k\)-schemes, if \(Y\) is smooth and proper over \(k\), exterior product induces an equivalence of stable \(\infty\)-categories
\[
\mathrm{DM}(X)^{\prime}
\otimes_{\mathrm{DM}(k)}
\mathrm{DM}(Y)^{\prime}
\xrightarrow{\ \boxtimes\ }
\mathrm{DM}(X\times_kY)^{\prime}.
\]
The quantification over \(X\), the coefficient convention, the meaning of the prime, and the meaning of this relative tensor product require confirmation from the complete mathematical setup. Thus this is a qualified transcription of the conjecture, not a fully specified standalone formulation.

\textbf{Supplied context.} Within its implicit setup, the report attributes to Jin--Yang the full faithfulness of
\[
\mathrm{DM}(X)\otimes_{\mathrm{DM}(k)}\mathrm{DM}(Y)
\longrightarrow \mathrm{DM}(X\times_kY).
\]
It describes the conjecture as concerning the construction of motivic local systems on the product from those on the factors, and mentions potential applications to the Langlands program. Neither the precise Jin--Yang theorem nor its hypotheses are supplied.

The following torsion statement is separate context. Let \(\ell\ne\operatorname{char}(k)\) be a prime. Write \(\mathrm{Mod}_{\Z/\ell}\) for the derived category of \(\Z/\ell\)-modules and \(D_{\mathrm{et}}(S;\Z/\ell)\) for the derived category of étale sheaves with these coefficients. The passage states a rigidity equivalence, written uniformly in the scheme variable as
\[
\mathrm{DM}(S)\otimes_{\mathrm{Mod}_{\Z}}\mathrm{Mod}_{\Z/\ell}
\simeq D_{\mathrm{et}}(S;\Z/\ell),
\]
and characterizes dualizable objects on the right as étale-locally constant objects locally given by perfect complexes of \(\Z/\ell\)-modules. For this torsion discussion it explicitly assumes that \(k\) is separably closed, so that
\(D_{\mathrm{et}}(\operatorname{Spec}k;\Z/\ell)\simeq\mathrm{Mod}_{\Z/\ell}\).
Under that assumption, its Theorem 2, attributed to \([\mathrm{HRS24}]\), states that if \(Y\) is smooth and either \(Y\) is proper or \(\operatorname{char}(k)=0\), exterior product gives
\[
D_{\mathrm{et}}(X;\Z/\ell)^{\prime}
\otimes_{\mathrm{Mod}_{\Z/\ell}}
D_{\mathrm{et}}(Y;\Z/\ell)^{\prime}
\simeq D_{\mathrm{et}}(X\times_kY;\Z/\ell)^{\prime}.
\]
The prime and the admissible class of schemes remain unspecified in this excerpt. The separably closed assumption is explicitly attached to the torsion discussion; it is not explicitly imposed on Conjecture 1.
\subsection{Short English statement}
The question asks whether, when \(Y\) is smooth and proper over \(k\), exterior product identifies the indicated categories of étale motives on \(X\) and \(Y\), tensored over motives of \(k\), with the indicated category on \(X\times_kY\). The coefficient convention, prime notation, hypotheses on \(X\) and \(k\), and categorical tensor product need clarification. The supplied dated curation labels the problem partially solved and says the full étale-motive conjecture remains open; this status is not independently verified.
\subsection{Sources}
\begin{itemize}
\item Supplied local UnsolvedMath registry record, permanent ID \(30006381\), exact source problem code \texttt{catalogue entry}, titled \emph{Categorical Künneth Formula for Étale Motives}. The record labels its statement \texttt{corrected\_verified}, with verification described as repairs only to formula delimiters, notation, grammar, and typography; this is supplied metadata, not independent mathematical verification. No proposer is identified.
\item Supplied report citation: \emph{Algebraic K-theory (2026)}, Oberwolfach Reports, pp.\ 1949--2012, DOI \texttt{10.4171/owr/2025/36}. \url{https://doi.org/10.4171/owr/2025/36}. The excerpt identifies the target as Conjecture 1, attributes an unprimed full-faithfulness theorem to Jin--Yang, and labels the torsion result Theorem 2, \([\mathrm{HRS24}]\). Complete bibliographic details for these attributions are not supplied.
\item Supplied literature source labeled ``Richarz--Scholbach primary paper'': \url{https://arxiv.org/abs/2503.14416}. The record does not supply its title or a theorem number identifying the relevant claims.
\item Supplied literature source labeled ``Oberwolfach report'': \url{https://doi.org/10.4171/owr/2025/36}.
\item Literature assessment checked at \texttt{2026-08-22}, copied from the supplied dated curation and not independently verified: the categorical Künneth equivalence is said to be proved for adic sheaves, étale torsion under the listed hypotheses, and cohomological motives, while the full étale-motive conjecture in the extracted statement is said to remain open. The supplied status is \texttt{partially\_solved}. The curation itself requests independent verification of time-sensitive status and cited claims.
\end{itemize}
\end{document}
